\begin{table}[t!]
\caption{
\textbf{Architecture details and model benchmarks for Chinese and code models.}
BPIC (defined in \cref{tab: main}) denotes the compression between the \main{} and outermost stage (bytes).
Each \model{} used (3,3)-DC, targeting an inner downsampling ratio of $9$.
However, the resulting BPIC was significantly different, indicating that code is much easier to compress than Chinese.
In terms of results, \model{} (2-stage) performs better than both \model{} (space) and BPE Transformer on Chinese, which is reflected in the downstreams.
On the other hand, \model{} (2-stage) achieves similar performance to \model{} (space) on code, and both \model{} models perform significantly better than Transformer.
}

\centering
\small

\begin{tabular}{lcccccccccccc}
\toprule
\multirow{2}{*}{\textsc{Model}}  & \multicolumn{4}{c}{\textsc{Chinese}}                                                  & & \multicolumn{3}{c}{\textsc{Code}} \\
\cmidrule{2-5}
\cmidrule{7-9}
                        & \textsc{BPIC}  & \textsc{Main arch.}            & \textsc{Val. BPB} $\downarrow$ & \textsc{XW-zh. acc.} $\uparrow$      &  & \textsc{BPIC} & \textsc{Main arch.}             & \textsc{Val. BPB} $\downarrow$ \\
\midrule
Transformer             & 3.62  & \colorbox{Main}{T15}  & \underline{0.7404}                & 0.599                       &  & 3.58 & \colorbox{Main}{T13}   & 0.3376                \\
\model{} (space)        & 3.38  & \colorbox{Main}{T19}  & 0.7478                & \underline{0.629}                       &  & 7.97 & \colorbox{Main}{T40}   & \underline{0.3163}                \\
\model{} (2-stage)      & 5.81  & \colorbox{Main}{T30}  & \textbf{0.7032}                & \textbf{0.663}                       &  & 7.23 & \colorbox{Main}{T28}   & \textbf{0.3161}                \\
\bottomrule
\end{tabular}
\label{tab: cd-supp}
\end{table}